The request-rooted model depends on a small set of invariants. The first invariant is that one \texttt{Request} should survive the lifecycle by default. Intake, routing, approval, execution, proof, and payout should remain attached to the same demand-side thread unless there is a real boundary of separate funding, separate ownership, separate market routing, or separate review. This prevents the system from scattering business truth across disconnected artifacts simply because the work became more detailed.

The second invariant is that internal decomposition does not imply a new business root. Sub-work created by a worker, planner, or runtime should become \texttt{FulfillmentStep} under an existing \texttt{Fulfillment}. This keeps execution decomposition subordinate to the accepted lane rather than letting every subtask spawn a new request or pseudo-order object. The rule matters because AI systems often create extra task objects by default, which makes lifecycle reasoning and payment accountability harder.

The third invariant is that commitment and fulfillment remain distinct. Selecting a route, accepting a proposal, or approving commercial terms is not the same as starting execution. Public or cross-actor work should preserve a \texttt{Commitment} gate before \texttt{Fulfillment} begins. Private owner-controlled shortcuts may exist in bounded cases, but those shortcuts should remain explicit exceptions rather than becoming the silent default.

The fourth invariant is that current summary state and full history are separated. The request root should store current truth such as the brief, visibility, high-level constraints, and active references. The append-only ledger and adjacent durable objects should store the full history of what actually happened. This protects both readability and auditability.

The fifth invariant is that completion must be evidenced according to execution mode. Digital work may close through document, code, or object references. Embodied work may require photographs, signatures, access logs, witness confirmation, or other field-specific proof. The system should therefore validate completion against proof obligations rather than against narrative closure alone.
